% =============================================================================
\section{Tarefa 1 --- Cinemática direta}
% =============================================================================

O Lite~6 é um manipulador série com seis juntas rotacionais: as três
primeiras (base, ombro e cotovelo) posicionam o pulso e as três últimas, com
eixos concorrentes, formam um pulso esférico que orienta a ferramenta. A
cinemática direta~\cite{sequeira2024,siciliano2009} dá a pose do flange
(posição $p$ e orientação $R$) em função de $q = (q_1,\dots,q_6)$.

\paragraph{Método e referenciais.} Usou-se a convenção de Denavit--Hartenberg
\emph{modificada}~\cite{craig2005}, a do documento do
fabricante~\cite{ufactory_dh} e a do controlador. Em cada elo $i$, $z_i$ está
sobre o eixo da junta $i$, $x_i$ sobre a normal comum a $z_i$ e $z_{i+1}$
(perpendicular aos dois, se se cruzam) e $y_i$ completa o referencial direto.
A base $\{0\}$ está no tampo da mesa, sobre o eixo da junta~1, com $z_0$
vertical (Figura~\ref{fig:frames}); $\{1\}$ e $\{2\}$ partilham a origem no
ombro e $\{4\}$ e $\{5\}$ no centro do pulso. Passa-se de $\{i-1\}$ para
$\{i\}$ com quatro movimentos elementares, e a pose do flange é o produto das
seis transformações, com os parâmetros da Tabela~\ref{tab:dh}:
\begin{equation}
\T{i-1}{i} = R_x(\alpha_{i-1})\,T_x(a_{i-1})\,R_z(\theta_i)\,T_z(d_i),
\ \ \theta_i = q_i + \text{offset}_i;
\qquad
\T{0}{6}(q) = \T{0}{1}\cdots\T{5}{6}
= \left[\begin{smallmatrix} R(q) & p(q)\\ 0 & 1\end{smallmatrix}\right].
\label{eq:fk}
\end{equation}

\begin{figure}[H]
\centering
\begin{minipage}[c]{0.44\textwidth}
\centering
\fig{Photos/referenciais.png}{\textwidth}
\captionof{figure}{Referenciais D-H com $q=0$, plano $x_0z_0$. Azul: $z_i$;
vermelho: $x_i$. $\{2\}$ e $\{5\}$ desenhados ao lado da origem, que é a de
$\{1\}$ e $\{4\}$. Cotas em mm.}
\label{fig:frames}
\end{minipage}\hfill
\begin{minipage}[c]{0.52\textwidth}
\centering
\small
\captionof{table}{Parâmetros D-H modificados do Lite~6~\cite{ufactory_dh}.}
\label{tab:dh}
\begin{tabular}{@{}ccccc@{}}
\toprule
\rowcolor{tabcab}
$i$ & offset $\theta_i$ & $d_i$ [mm] & $\alpha_{i-1}$ & $a_{i-1}$ [mm] \\
\midrule
1 & \ang{0}   & 243,3 & \ang{0}   & 0 \\
2 & \ang{-90} & 0     & \ang{-90} & 0 \\
3 & \ang{-90} & 0     & \ang{180} & 200 \\
4 & \ang{0}   & 227,6 & \ang{90}  & 87 \\
5 & \ang{0}   & 0     & \ang{90}  & 0 \\
6 & \ang{0}   & 61,5  & \ang{-90} & 0 \\
\bottomrule
\end{tabular}
\end{minipage}
\end{figure}

Como o pulso é esférico, a posição separa-se da orientação: o centro do pulso
só depende das três primeiras juntas,
\begin{equation}
p_w = \begin{bmatrix} r\cos q_1\\ r\sin q_1\\ h \end{bmatrix},\qquad
\begin{aligned}
r &= 200\sin q_2 + 87\cos(q_3-q_2) + 227{,}6\,\sin(q_3-q_2),\\
h &= 243{,}3 + 200\cos q_2 + 87\sin(q_3-q_2) - 227{,}6\,\cos(q_3-q_2),
\end{aligned}
\label{eq:pulso}
\end{equation}
e o flange fica \SI{61.5}{\milli\meter} à frente, ao longo do eixo da
ferramenta: $p = p_w + 61{,}5\,R\,\hat e_z$. A orientação exprime-se em
\emph{roll}, \emph{pitch} e \emph{yaw},
$R = R_z(\text{yaw})\,R_y(\text{pitch})\,R_x(\text{roll})$, como no
SDK~\cite{xarm_sdk}. A ponta da pinça está a $L = \SI{83.6}{\milli\meter}$
ao longo de $z_6$: $\T{0}{\text{ponta}} = \T{0}{6}\,T_z(L)$.

\paragraph{Verificação.} \textbf{(i)} Com $q=0$, \eqref{eq:pulso} dá
$p = (87;\,0;\,154{,}2)$~mm com \emph{roll} de \ang{180}, como na
Figura~\ref{fig:frames}. \textbf{(ii)} A tabela D-H clássica do mesmo
documento, implementada à parte, coincide em \num{10000} configurações
aleatórias a menos de \SI{2e-13}{\milli\meter}. \textbf{(iii)} O modelo 3D
do fabricante coincide a menos de \SI{0.7}{\milli\meter} no flange ao longo
do swing (Secção~\ref{sec:contacto}). \textbf{(iv)} Face ao controlador, nos
100 pontos da Tarefa~4, a diferença é de \SI{1.2}{\milli\meter} em média
(máximo \SI{2.0}{\milli\meter}) e \ang{0.08}; com $q=0$ o controlador indica
o flange em $(86{,}82;\,-0{,}61;\,154{,}32)$~mm, a \SI{0.6}{\milli\meter} do
nosso, porque aplica uma calibração própria (Secção~\ref{sec:ik}).
